\documentclass[10pt,a4paper]{amsart}
\usepackage[margin=19mm]{geometry}
\usepackage{amsmath,amssymb,mathtools}
\usepackage[scr=rsfs,cal=euler]{mathalfa}
\usepackage{xfrac}
\usepackage[unicode,hidelinks,bookmarks=false]{hyperref}
\newcommand{\ie}{i.e.\ }
\newcommand{\cf}{cf.\ }
\newcommand{\resp}{respectively\ }
\newcommand{\Subsection}{\subsection}
\providecommand{\nobreakdash}{\nobreak\hbox{-}}
\setlength{\emergencystretch}{2em}
\multlinegap0pt
\usepackage{bm}
\usepackage{bbm}
\usepackage{dsfont}
\usepackage{xspace}
\usepackage{cancel}
\usepackage{mathtools}
\usepackage{amsthm}
\makeatletter
\@ifundefined{theorem}{\newtheorem{theorem}{Theorem}}{}
\@ifundefined{lemma}{\newtheorem{lemma}[theorem]{Lemma}}{}
\makeatother
\newcommand{\abs}[1]{{\left|#1\right|}}
\title[Local stability for a class of Saint-Venant type inequalitie]{Local stability for a class of Saint-Venant type inequalities}
\author{Luca Barbato, Francesco Salerno}
\date{Source: arXiv:2606.27031v1 (CC BY 4.0)}
\begin{document}
\maketitle

\noindent\emph{Source and licence.} Local stability for a class of Saint-Venant type inequalities. Version arXiv:2606.27031v1 (CC BY 4.0), distributed under the Creative Commons Attribution 4.0 International licence, as recorded in the arXiv licence metadata for this exact version and shown on the abstract page https://arxiv.org/abs/2606.27031v1. Full rights notice: licenses/arxiv-2606.27031-CC-BY-4.0.txt.\par\smallskip

\noindent\textit{Editorial scope.} Counted unit: the Lemma whose source label is \texttt{dersec(3)} in the article (source0.tex lines 401--410, source label \texttt{dersec(3)}), delivered together with the article's own complete original proof of it (source0.tex lines 411--482, one \texttt{proof} environment of 72 lines). No mathematics is written, repaired or re-derived: statement and proof are the article's own text line for line. Required context retained uncounted: Problema(u) (lines 124--129); lemmaA.1 (lines 159--181); Problema(u.) (lines 374--379); Problema(p) (lines 381--386). Every definition, display and label of those lines is retained unchanged. Omitted from this package and not redistributed: 1--123, 130--158, 182--373, 380--380, 387--400, 483--1001. Nothing inside a retained excerpt is omitted or altered: every retained body is delivered exactly as the article's lines are written, so no omission receipt of any kind accompanies this package. Citations: the retained excerpts carry no citation command at all, so no bibliography entry of the article is transcribed and no reference is redistributed. Reference adaptation: none was needed and none was made; every reference inside the delivered unit resolves inside it, so no unresolved reference or citation is produced. Printed slips kept literal: no printed word, symbol, hypothesis or number of the counted statement or of its proof was altered. Layout and class: the article's own class is \texttt{article} with packages including geometry, fontenc, inputenc, babel, amssymb, lipsum, lmodern, amsmath, mathrsfs, cancel, amsthm, bm, mathtools, braket; the wrapper is the lane's pinned \texttt{amsart} editorial wrapper at 10pt on A4, which restates the theorem-like environments the retained text needs and adds the article's own macro definitions that the retained text needs (\texttt{\textbackslash abs}), with extra wrapper packages (bm, bbm, dsfont, xspace, cancel, mathtools). The delivered numbering is the unit's own. Assets: no retained span contains a figure environment, a graphics-inclusion command, a PDF-page inclusion, a table or a TikZ picture; no image file is copied into this package, the package declares no asset dependency and no image file is redistributed with it. Licence: version arXiv:2606.27031v1 of this article is distributed under the Creative Commons Attribution 4.0 International licence, as recorded in the arXiv licence metadata for this exact version and shown on the abstract page https://arxiv.org/abs/2606.27031v1. A full rights notice is delivered at licenses/arxiv-2606.27031-CC-BY-4.0.txt. Characters: every retained excerpt is pure ASCII, so the delivered engine, XeLaTeX with its default Latin Modern text font, reports no missing character.
\begin{equation}\label{Problema(u)}
    \begin{cases}
        -\Delta u=1 & \text{in } \Omega,\\
        u=0 & \text{on } \partial\Omega.
    \end{cases}
\end{equation}

\begin{lemma}\label{lemmaA.1}
    Given $\gamma\in(0,1]$ there exists $\delta_e=\delta_e(n,\gamma)>0$ and a modulus of continuity $\omega_e$ such that, for every nearly spherical set $\Omega$ parametrized by $\varphi$ with $\|\varphi\|_{C^{2,\gamma}(\partial B_1)}\leq\delta_e$ and $\abs{\Omega}=\abs{B_1}$, we can find an autonomous vector field $V_\varphi$ for which the following holds true:
    \begin{itemize}
        \item [(i)] $\mathrm{div}\,V_\varphi=0$ in a $\delta_e$-neighborhood of $\partial B_1$,
        \item [(ii)] if $\Phi_t:=\Phi(t,x)$ is the flow of $V_\varphi$, i.e.
        \begin{equation*}
            \partial_t \Phi_t=V_\varphi(\Phi_t),\quad \Phi_0(x)=x,
        \end{equation*}
        then $\Phi_1(\partial B_1)=\partial\Omega$ and $\abs{\Phi_t(B_1)}=\abs{B_1}$ for all $t\in[0,1]$,
        \item [(iii)] We have  
        \begin{equation*}
            \|\Phi_t-Id\|_{C^{2,\gamma}}\leq\omega_e(\|\varphi\|_{C^{2,\gamma}(\partial B_1)})\quad \text{for every }t\in[0,1],
        \end{equation*}
        \begin{equation*}
            \|\varphi-(V_\varphi\cdot\nu_0)\|_{H^{1/2}(\partial B_1)}\leq\omega_e(\|\varphi\|_{L^\infty(\partial B_1)})\|\varphi\|_{H^{1/2}(\partial B_1)},
        \end{equation*}
        and
        \begin{equation*}
            (V_\varphi\cdot\theta)\circ\Phi_t-V_\varphi\cdot\nu_0=(V_\varphi\cdot\nu_0)\psi_t \quad \text{on }\partial B_1,
        \end{equation*}
        with $\|\psi_t\|_{C^{2,\gamma}(\partial B_1)}\leq\omega_e(\|\varphi\|_{C^{2,\gamma}(\partial B_1)})$.
    \end{itemize}
\end{lemma}

\begin{equation}\label{Problema(u.)}
    \begin{cases}
        -\Delta \dot{u}_t=0 & \text{ in }\Omega(t)\\
        \dot{u}_t=-\nabla u_t\cdot V & \text{ on } \partial\Omega(t).
    \end{cases}
\end{equation}

\begin{equation}\label{Problema(p)}
    \begin{cases}
        -\Delta p_t= j'(u_t) & \text{ in }\Omega(t)\\
        p_t=0 & \text{ on } \partial\Omega(t).
    \end{cases}
\end{equation}

\begin{lemma}[Second Variation]
    Under the assumptions of Lemma \ref{lemmaA.1}, we have
    \begin{equation}\label{dersec(3)}
    \begin{split}
        f''(t)=&-2\int_{\Omega(t)}\nabla \dot{u}_t\cdot\nabla\dot{p}_t \,dx+\int_{\Omega(t)}j''(u_t)\dot{u}_t^2\,dx+\int_{\partial\Omega(t)}\frac{\abs{\nabla p_t}}{\abs{\nabla u_t}}\langle D^2u_t.\nabla u_t,V^\tau\rangle(V\cdot\nu_t)\,d\mathcal{H}^{n-1}\\
        &+\int_{\partial\Omega(t)}\frac{\abs{\nabla u_t}}{\abs{\nabla p_t}}\langle D^2p_t.\nabla p_t,V^\tau\rangle(V\cdot\nu_t)\,d\mathcal{H}^{n-1}+\int_{\partial\Omega(t)}\abs{\nabla p_t}(V\cdot\nu_t)^2\,d\mathcal{H}^{n-1}\\
        &+\int_{\partial\Omega(t)}\abs{\nabla u_t}j'(u_t)(V\cdot\nu_t)^2\,d\mathcal{H}^{n-1}-2\int_{\partial\Omega(t)}\abs{\nabla u_t}\abs{\nabla p_t}\mathscr{H}_{\partial\Omega(t)}(V\cdot\nu_t)^2\,d\mathcal{H}^{n-1}.
    \end{split}
\end{equation}
\end{lemma}

\begin{proof}
  Since we want to apply the Hadamard's formula to $f'(t)$, we rewrite the boundary integral, using the divergence Theorem, as follows
\begin{equation*}
    f'(t)=\int_{\Omega(t)}\mathrm{div}(\abs{\nabla u_t}\abs{\nabla p_t}V)\,dx.
\end{equation*}
Then we can compute the second derivative
\begin{equation}\label{dersec(1)}
    f''(t)=\int_{\Omega(t)}\frac{\partial}{\partial t}\mathrm{div}(\abs{\nabla u_t}\abs{\nabla p_t}V)\,dx+\int_{\partial\Omega(t)}\mathrm{div}(\abs{\nabla u_t}\abs{\nabla p_t}V)(V\cdot\nu_t)\,d\mathcal{H}^{n-1}.
\end{equation}
Regarding the first integral, since $V$ is an autonomous vector field, applying again the divergence Theorem, we have that
\begin{equation}\label{parte1}
    \begin{split}
        \int_{\Omega(t)}\frac{\partial}{\partial t}\mathrm{div}(\abs{\nabla u_t}\abs{\nabla p_t}V)\,dx&=\int_{\Omega(t)}\mathrm{div}\left(\frac{\partial}{\partial t}(\abs{\nabla u_t}\abs{\nabla p_t})V\right)\,dx=\int_{\partial\Omega(t)}\frac{\partial}{\partial t}(\abs{\nabla u_t}\abs{\nabla p_t})(V\cdot\nu_t)\,d\mathcal{H}^{n-1}\\
        &=\int_{\partial\Omega(t)}\left[\abs{\nabla p_t}\frac{\nabla u_t}{\abs{\nabla u_t}}\cdot\nabla \dot{u}_t+\abs{\nabla u_t}\frac{\nabla p_t}{\abs{\nabla p_t}}\cdot\nabla \dot{p}_t \right](V\cdot\nu_t)\,d\mathcal{H}^{n-1}\\
        &=-\int_{\partial\Omega(t)}\left[\abs{\nabla p_t}\frac{\partial\dot{u}_t}{\partial\nu_t}+\abs{\nabla u_t}\frac{\partial\dot{p}_t}{\partial\nu_t} \right](V\cdot\nu_t)\,d\mathcal{H}^{n-1}.
    \end{split}
\end{equation}
Now we want to rewrite the second term in \eqref{dersec(1)}. For simplicity we recall that
\begin{equation*}
    \nabla(\abs{\nabla g})\cdot V=\frac{\langle D^2g.\nabla g, V\rangle}{\abs{\nabla g}}.
\end{equation*}
Then, recalling again that $\mathrm{div}V=0$ in a neighborhood of $\partial B_1$, the boundary integral in \eqref{dersec(1)} can be rewritten as follows
\begin{equation}\label{parte2}
    \begin{split}
        \int_{\partial\Omega(t)}\mathrm{div}(\abs{\nabla u_t}\abs{\nabla p_t}V)(V\cdot\nu_t)\,d\mathcal{H}^{n-1}&=\int_{\partial\Omega(t)}\abs{\nabla p_t}\frac{\langle D^2u_t.\nabla u_t, V\rangle}{\abs{\nabla u_t}}(V\cdot\nu_t)\,d\mathcal{H}^{n-1}\\
        &+\int_{\partial\Omega(t)}\abs{\nabla u_t}\frac{\langle D^2p_t.\nabla p_t, V\rangle}{\abs{\nabla p_t}}(V\cdot\nu_t)\,d\mathcal{H}^{n-1}.
    \end{split}
\end{equation}
Finally, combining \eqref{parte1}-\eqref{parte2}, we have that
\begin{equation*}
    \begin{split}
        f''(t)=&-\int_{\partial\Omega(t)}\left[\abs{\nabla p_t}\frac{\partial\dot{u}_t}{\partial\nu_t}+\abs{\nabla u_t}\frac{\partial\dot{p}_t}{\partial\nu_t} \right](V\cdot\nu_t)\,d\mathcal{H}^{n-1}+\int_{\partial\Omega(t)}\abs{\nabla p_t}\frac{\langle D^2u_t.\nabla u_t, V\rangle}{\abs{\nabla u_t}}(V\cdot\nu_t)\,d\mathcal{H}^{n-1}\\
        &+\int_{\partial\Omega(t)}\abs{\nabla u_t}\frac{\langle D^2p_t.\nabla p_t, V\rangle}{\abs{\nabla p_t}}(V\cdot\nu_t)\,d\mathcal{H}^{n-1}.
    \end{split}
\end{equation*}
We remark that, from \eqref{Problema(p)}, $\dot{p}_t$ satisfies
\begin{equation}\label{Problema(p.)}
    \begin{cases}
        -\Delta \dot{p}_t=j''(u_t)\dot{u}_t & \text{ in }\Omega(t)\\
        \dot{p}_t=-\nabla p_t\cdot V & \text{ on } \partial\Omega(t).
    \end{cases}
\end{equation}
Now we decompose $V=V^\tau+(V\cdot\nu_t)\nu_t$, and, using the boundary condition in \eqref{Problema(u.)}-\eqref{Problema(p.)}, we find the following
\begin{equation}\label{dersec(2)}
    \begin{split}
        f''(t)=&-\int_{\partial\Omega(t)}\left[\dot{p}_t\frac{\partial\dot{u}_t}{\partial\nu_t}+\dot{u}_t\frac{\partial\dot{p}_t}{\partial\nu_t} \right]\,d\mathcal{H}^{n-1}-\int_{\partial\Omega(t)}\abs{\nabla p_t}\langle D^2u_t.\nu_t,V^\tau\rangle(V\cdot\nu_t)\,d\mathcal{H}^{n-1}\\
        &-\int_{\partial\Omega(t)}\abs{\nabla p_t}\langle D^2u_t.\nu_t,\nu_t\rangle(V\cdot\nu_t)^2\,d\mathcal{H}^{n-1}-\int_{\partial\Omega(t)}\abs{\nabla u_t}\langle D^2p_t.\nu_t,V^\tau\rangle(V\cdot\nu_t)\,d\mathcal{H}^{n-1}\\
        &-\int_{\partial\Omega(t)}\abs{\nabla u_t}\langle D^2p_t.\nu_t,\nu_t\rangle(V\cdot\nu_t)^2\,d\mathcal{H}^{n-1}.
    \end{split}
\end{equation}
We remark that, whenever $g$ is a $C^2$ concave functions that vanish on the boundary of $\Omega(t)$, we have
\begin{equation}\label{infLapl}
    \langle D^2g.\nu_t,\nu_t\rangle=\Delta g+\abs{\nabla g}\mathscr{H}_{\partial\Omega(t)},
\end{equation}
where $\mathscr{H}_{\partial\Omega(t)}$ is the mean curvature of the boundary of $\Omega(t)$. Then, from \eqref{Problema(u)}-\eqref{Problema(p)}-\eqref{dersec(2)}-\eqref{infLapl}, we find
\begin{equation*}
    \begin{split}
        f''(t)=&-\int_{\partial\Omega(t)}\left[\dot{p}_t\frac{\partial\dot{u}_t}{\partial\nu_t}+\dot{u}_t\frac{\partial\dot{p}_t}{\partial\nu_t} \right]\,d\mathcal{H}^{n-1}-\int_{\partial\Omega(t)}\abs{\nabla p_t}\langle D^2u_t.\nu_t,V^\tau\rangle(V\cdot\nu_t)\,d\mathcal{H}^{n-1}\\
        &-\int_{\partial\Omega(t)}\abs{\nabla u_t}\langle D^2p_t.\nu_t,V^\tau\rangle(V\cdot\nu_t)\,d\mathcal{H}^{n-1}+\int_{\partial\Omega(t)}\abs{\nabla p_t}(V\cdot\nu_t)^2\,d\mathcal{H}^{n-1}\\
        &+\int_{\partial\Omega(t)}\abs{\nabla u_t}j'(u_t)(V\cdot\nu_t)^2\,d\mathcal{H}^{n-1}-2\int_{\partial\Omega(t)}\abs{\nabla u_t}\abs{\nabla p_t}\mathscr{H}_{\partial\Omega(t)}(V\cdot\nu_t)^2\,d\mathcal{H}^{n-1}.
    \end{split}
\end{equation*}

Finally, using the divergence Theorem on the first integral, since $u_t=p_t=0$ on $\partial\Omega(t)$, from \eqref{Problema(u.)}-\eqref{Problema(p.)}, we can rewrite $f''(t)$ as follows
\begin{equation*}
    \begin{split}
        f''(t)=&-2\int_{\Omega(t)}\nabla \dot{u}_t\cdot\nabla\dot{p}_t \,dx+\int_{\Omega(t)}j''(u_t)\dot{u}_t^2\,dx+\int_{\partial\Omega(t)}\frac{\abs{\nabla p_t}}{\abs{\nabla u_t}}\langle D^2u_t.\nabla u_t,V^\tau\rangle(V\cdot\nu_t)\,d\mathcal{H}^{n-1}\\
        &+\int_{\partial\Omega(t)}\frac{\abs{\nabla u_t}}{\abs{\nabla p_t}}\langle D^2p_t.\nabla p_t,V^\tau\rangle(V\cdot\nu_t)\,d\mathcal{H}^{n-1}+\int_{\partial\Omega(t)}\abs{\nabla p_t}(V\cdot\nu_t)^2\,d\mathcal{H}^{n-1}\\
        &+\int_{\partial\Omega(t)}\abs{\nabla u_t}j'(u_t)(V\cdot\nu_t)^2\,d\mathcal{H}^{n-1}-2\int_{\partial\Omega(t)}\abs{\nabla u_t}\abs{\nabla p_t}\mathscr{H}_{\partial\Omega(t)}(V\cdot\nu_t)^2\,d\mathcal{H}^{n-1}.
    \end{split}
\end{equation*}
\end{proof}

\end{document}
